% =====================================================================
%  Lecture 11 -- Home Subscriber Server (HSS)             (English edition)
% =====================================================================
\documentclass[11pt,a4paper]{article}
\providecommand\ACNroot{../..}
\input{\ACNroot/templates/lecture-en}

\ACNset{number}{11}
\ACNset{title}{Home Subscriber Server (HSS)}
\ACNset{shorttitle}{HSS}
\ACNset{titleother}{کارساز مشترکین خانگی (HSS)}
\ACNset{part}{Part II -- 4G access and core (Lectures 5--13)}
\ACNset{time}{6--8 hours (3 h reading, 3--5 h exercises)}
\ACNset{level}{Intermediate}
\ACNset{prereq}{Lectures 9 and 10; basic probability; symmetric cryptography}
\ACNset{gen}{4}

\begin{document}
\ACNcover

% ---------------------------------------------------------------------
\section{Learning objectives}
After studying this lecture you should be able to:
\begin{objectives}
  \item \textbf{Describe} the HSS as the union of HLR, AuC and IMS user-database functions, and
        \textbf{list} its interfaces.
  \item \textbf{Explain} the subscription data model of the EPS and how the MME uses it.
  \item \textbf{Trace} authentication-vector generation in the AuC and \textbf{explain} sequence-number
        management and re-synchronisation.
  \item \textbf{Size} HSS storage and vector-request load, and \textbf{evaluate} the trade-off of
        requesting several vectors at once.
  \item \textbf{Compute} the availability of geo-redundant HSS deployments.
  \item \textbf{Relate} the HSS to the 5G UDM/AUSF/UDR and to data-layered (UDC) architectures.
\end{objectives}

% ---------------------------------------------------------------------
\section{Prerequisite knowledge}
This lecture builds on the EPS-AKA procedure and key hierarchy of Lecture~9 and on the S6a application of
Lecture~10. You should know what a block cipher and a message authentication code are, and be able to
compute probabilities of independent events (for the availability models).

\paragraph{Key terms} The vocabulary of this lecture with its Persian equivalents:
\begin{keyterms}
HSS & \fa{کارساز مشترکین خانگی} & Master subscriber database of a 3GPP network (HLR + AuC + IMS data). \\
HLR & \fa{ثبت‌کنندهٔ موقعیت خانگی} & 2G/3G subscriber database reached over MAP. \\
AuC & \fa{مرکز احراز هویت} & Stores the secret key $K$ and generates authentication vectors. \\
IMSI & \fa{شناسهٔ بین‌المللی مشترک موبایل} & Permanent subscriber identity (MCC + MNC + MSIN). \\
Authentication vector & \fa{بردار احراز هویت} & (RAND, XRES, AUTN, $K_{\mathrm{ASME}}$) for one EPS-AKA run. \\
SQN & \fa{شمارهٔ ترتیب} & 48-bit counter that proves the freshness of a vector. \\
AUTS & \fa{نشانهٔ همگام‌سازی مجدد} & Token the USIM returns to re-synchronise SQN. \\
OPc & \fa{کلید پیکربندی اپراتور} & Operator-variant value derived from OP and $K$ for MILENAGE. \\
UDC / UDR & \fa{همگرایی دادهٔ کاربر / مخزن} & Architecture separating application logic from a shared data repository. \\
Geo-redundancy & \fa{افزونگی جغرافیایی} & Replicating a function across distant sites. \\
\end{keyterms}

% ---------------------------------------------------------------------
\section{Introduction and motivation}
A mobile network is, among other things, a very large and very available database. Every subscriber
has a permanent identity, a secret key shared only with the SIM, a profile of allowed services and a
current location; every attach, many handovers and every incoming call read or write some of these.
In 2G and 3G these data lived in the Home Location Register (HLR) and the Authentication Centre (AuC);
with IMS (Release~5) 3GPP introduced the \emph{Home Subscriber Server} (HSS), which contains the HLR
and AuC functions and adds the IMS user data \cite{ts23002}. In the EPS the HSS is the single
subscriber authority for the MME, reached over Diameter S6a (Lecture~10).

The HSS is rarely discussed in textbooks because it does not move packets, yet an HSS outage stops
every new attach in a country. Its design questions -- how to keep $10^7$--$10^8$ profiles consistent
across sites, how to generate fresh authentication vectors at thousands per second, how to recover
sequence numbers after a failure -- are typical of carrier-grade distributed systems.

% ---------------------------------------------------------------------
\section{Fundamental concepts}

\subsection{Functions of the HSS}
\begin{itemize}
  \item \textbf{Identification and addressing}: IMSI, MSISDN, IMS public and private identities
        \cite{ts23003}.
  \item \textbf{Authentication (AuC)}: stores $K$ (and OPc) and generates authentication vectors
        \cite{ts33102,ts33401}.
  \item \textbf{Subscription management}: stores the profile -- allowed APNs, QoS, roaming and access
        restrictions -- and pushes changes to the serving node.
  \item \textbf{Mobility management}: records the serving MME/SGSN/MSC so that the network can reach the
        subscriber, and cancels the old registration on a change of serving node.
  \item \textbf{IMS support}: stores the serving S-CSCF, IMS service profiles and initial filter
        criteria, and answers application servers \cite{ts23228}.
\end{itemize}

\subsection{Interfaces}
\begin{table}[htbp]
\centering\small
\caption{Main interfaces of the HSS.}\label{tab:ifaces}
\rowcolors{2}{acnmist}{white}
\begin{tabularx}{\linewidth}{@{}l l l Y@{}}
\toprule
\rowcolor{white}Interface & Peer & Protocol & Use \\
\midrule
S6a / S6d & MME / SGSN (EPS)   & Diameter & Vectors, location update, subscription data \cite{ts29272} \\
Gr        & SGSN (2G/3G)       & MAP      & Same functions for GPRS/UMTS \cite{ts29002} \\
C / D     & GMSC / MSC-VLR     & MAP      & Call routing information, CS location update \cite{ts29002} \\
Cx        & I-CSCF / S-CSCF    & Diameter & IMS registration, S-CSCF assignment \\
Sh        & Application server & Diameter & Transparent service data, notifications \\
SWx       & 3GPP AAA server    & Diameter & Authentication for non-3GPP access (e.g.\ Wi-Fi) \\
\bottomrule
\end{tabularx}
\end{table}

\subsection{Subscription data}
\begin{definitionbox}{EPS subscription profile}
The profile that the HSS sends in ULA (and later in IDR) contains, among others, the subscriber status,
MSISDN, network-access mode, operator-determined barring, the subscribed UE-AMBR, access-restriction
data (e.g.\ ``no E-UTRAN''), the periodic TAU timer and an \emph{APN-configuration profile}: a list of
APN configurations, each with a context identifier, the APN name (Service-Selection), PDN type, the
default bearer's QCI and ARP and the APN-AMBR \cite{ts29272}.
\end{definitionbox}
The MME caches the profile while the subscriber is registered with it. The HSS remains the master: it
changes the cached copy with Insert-Subscriber-Data (IDR) or Delete-Subscriber-Data (DSR), and
withdraws it with Cancel-Location (CLR).

% ---------------------------------------------------------------------
\section{Detailed technical explanation}

\subsection{Authentication-vector generation}
For each vector the AuC chooses a random RAND, takes the next sequence number SQN and evaluates five
functions of $K$, RAND, SQN and AMF (\cref{fig:av}): $f_1$ gives the network authentication code MAC,
$f_2$ the expected response XRES, $f_3$ and $f_4$ the cipher and integrity keys CK and IK, and $f_5$ the
anonymity key AK that conceals SQN \cite{ts33102}. The recommended algorithm set, MILENAGE, builds
all five from one 128-bit block cipher (AES) and the operator constant OPc \cite{ts35206}. For E-UTRAN the
HSS sets the \emph{AMF separation bit} to 1 and derives $K_{\mathrm{ASME}}$ from CK and IK, binding the
vector to the serving network \cite{ts33401}. The vector is
\begin{equation}\label{eq:autn}
  \mathrm{AV} = \bigl(\mathrm{RAND},\ \mathrm{XRES},\ \mathrm{AUTN},\ K_{\mathrm{ASME}}\bigr),\qquad
  \mathrm{AUTN} = (\mathrm{SQN}\xor\mathrm{AK})\concat\mathrm{AMF}\concat\mathrm{MAC}.
\end{equation}
The secret $K$ never leaves the AuC; the MME receives only vectors.

\begin{acnfigure}
  \resizebox{0.92\linewidth}{!}{\acnfig{lecture-11/av-generation}}
  \caption{Authentication-vector generation in the AuC with the functions $f_1$--$f_5$.}
  \label{fig:av}
\end{acnfigure}

\subsection{Sequence numbers and freshness}
The USIM accepts AUTN only if MAC is correct \emph{and} SQN is fresh. TS~33.102 describes, in an
informative annex, a scheme in which SQN is split into a 43-bit counter SEQ and a 5-bit index IND
($\mathrm{SQN} = \mathrm{SEQ}\concat\mathrm{IND}$) \cite{ts33102}. The USIM keeps an array of the highest SEQ
accepted for each of the $2^5 = 32$ values of IND, and accepts a new SQN if its SEQ exceeds the stored
value in the slot IND. The array lets vectors that were generated in one order be used in another --
for example when two MMEs hold batches of vectors for the same subscriber -- without spurious
failures. An optional age check also rejects SQN values that are too far behind the highest accepted
one.

If the SQN is not fresh, the USIM answers with a \emph{synchronisation failure} carrying AUTS, which
contains its own SQN concealed with a key derived like AK and protected by a MAC (MAC-S). The MME sends
RAND and AUTS to the HSS in a new AIR (Re-Synchronization-Info); the AuC verifies MAC-S, sets its
counter to the USIM's value and returns fresh vectors (\cref{fig:resync}).

\subsection{Mobility management in the HSS}
On each ULR the HSS records the serving MME (Origin-Host and realm). If another MME was registered, the
HSS sends it CLR so that it deletes its context; this guarantees one serving node per subscriber. An MME
that deletes a long-inactive subscriber sends Purge-UE (PUR), after which the HSS marks the subscriber as
purged. After a restart, an HSS that may have lost registrations can send Reset (RSR) to the MMEs, which
then refresh the affected subscribers at their next contact \cite{ts29272}.

\subsection{Implementation: front-ends, data store and UDC}
Commercial HSSs separate stateless \emph{front-ends} that terminate MAP and Diameter and run the
application logic from a replicated, mostly in-memory \emph{data store}. 3GPP generalised this as
\emph{User Data Convergence} (UDC): HLR/HSS, AuC and other front-ends share one User Data Repository
(UDR) over the Ud interface, so that a subscriber has one record instead of a copy per node
(\cref{fig:arch}). The 5G core keeps the same split: UDM and AUSF are front-ends, UDR is the repository
(Lectures~18--19). Because the HSS is a single point of control, the data store is replicated across
geographically separated sites.

\begin{acnfigure}
  \resizebox{\linewidth}{!}{\acnfig{lecture-11/hss-architecture}}
  \caption{Functional view of the HSS with its interfaces and the front-end/data-store split.}
  \label{fig:arch}
\end{acnfigure}

% ---------------------------------------------------------------------
\section{Mathematical models}

\subsection{Sequence-number lifetime}\label{sec:sqn}
\paragraph{Model} With a $b$-bit SEQ and one authentication every $\tau$ seconds, the counter lasts
\begin{equation}\label{eq:sqnlife}
  T_{\mathrm{life}} = 2^{b}\,\tau .
\end{equation}
\paragraph{Example} $b = 43$, one authentication per minute: $T_{\mathrm{life}} = 8.8\times10^{12}$\,min
$\approx 1.7\times10^{7}$ years. Exhaustion is not a concern; the counter is large so that schemes that
skip values (for example a SEQ derived from a clock) remain possible.

\subsection{Storage}\label{sec:storage}
\paragraph{Model} $N$ subscribers with a mean profile size $s$, replicated $R$ times, with a fill factor
$\eta$ (fraction of memory that may be used):
\begin{equation}\label{eq:storage}
  M = \frac{N\,s\,R}{\eta}.
\end{equation}
\paragraph{Example} $N = 10^7$, $s = 4$\,kB, $R = 3$, $\eta = 1$: $M = 120$\,GB -- small enough to keep
in memory on a few servers. Storage is therefore not the hard problem; consistency and availability
are.

\subsection{Vector requests}\label{sec:air}
\paragraph{Model} If authentications occur at rate $\lambda_{\mathrm{auth}}$ and the MME requests $n$
vectors per AIR and uses all of them, the AIR rate is
\begin{equation}\label{eq:air}
  \lambda_{\mathrm{AIR}} = \frac{\lambda_{\mathrm{auth}}}{n}.
\end{equation}
\paragraph{Example} $10^7$ subscribers authenticated 10 times a day: $\lambda_{\mathrm{auth}} = 1157$\,s$^{-1}$;
$n = 5$ gives $\lambda_{\mathrm{AIR}} = 231$\,s$^{-1}$. \paragraph{Limitation} Vectors left unused when the
subscriber moves to another MME are wasted, and stored vectors consume MME memory, so $n$ is usually
small (\cref{exr:batch}).

\subsection{Availability of redundant sites}\label{sec:avail}
\paragraph{Model} Each site is available with probability $a$ (unavailability $u = 1-a$), independently
of the others. If service needs at least $k$ of $N$ sites,
\begin{equation}\label{eq:avail}
  A_{k\text{-of-}N} = \sum_{i=k}^{N}\binom{N}{i}a^{i}(1-a)^{N-i}.
\end{equation}
Special cases: $A_{1\text{-of-}2} = 1-u^2$; $A_{2\text{-of-}3} = 1-(3u^2-2u^3)$; $A_{2\text{-of-}2} = a^2$.
\paragraph{Example} $a = 0.999$ (8.76\,h downtime per year): two sites each able to carry the full load
give $u^2 = 10^{-6}$, i.e.\ 31.5\,s per year; two sites that are both needed give 17.5\,h per year
(\cref{fig:avail}).
\paragraph{Assumptions and interpretation} Independence is the critical assumption: software bugs,
configuration errors and replication faults hit all sites at once, and in practice dominate HSS outages.
The model gives an upper bound on what redundancy can achieve, not a prediction.

\begin{acnfigure}
  \acnfig{lecture-11/availability}
  \caption{Annual downtime as a function of single-site unavailability for four deployment options
           (independent failures assumed).}
  \label{fig:avail}
\end{acnfigure}

% ---------------------------------------------------------------------
\section{Network architecture and diagrams}
\Cref{fig:arch} shows the HSS between the circuit and packet core of 2G/3G, the EPS and the IMS. The
same subscriber record serves all of them: a VoLTE subscriber is authenticated for E-UTRAN over S6a,
registered in the IMS over Cx, and reached by an MSC over MAP if the call falls back to the circuit
domain. In the EPS, the HSS has exactly one control-plane peer type for mobility, the MME, and is
usually reached via DRAs (Lecture~10).

% ---------------------------------------------------------------------
\section{Step-by-step procedures}

\begin{procedurebox}{Generating an E-UTRAN vector in the AuC}\label{proc:av}
\begin{steps}
  \item Receive AIR with the IMSI and the serving-network identity (Visited-PLMN-Id).
  \item Look up $K$ and OPc; generate a fresh 128-bit RAND; take the next SQN (SEQ and IND).
  \item Compute MAC $= f_1$, XRES $= f_2$, CK $= f_3$, IK $= f_4$, AK $= f_5$ with AMF separation bit = 1.
  \item Form AUTN with \cref{eq:autn}; derive $K_{\mathrm{ASME}}$ from CK, IK, the serving-network identity
        and SQN$\xor$AK.
  \item Return the vectors in AIA; the HSS discards CK and IK.
\end{steps}
\end{procedurebox}

\begin{procedurebox}{Sequence-number re-synchronisation}\label{proc:resync}
\begin{steps}
  \item The USIM verifies MAC, recovers SQN and finds it not fresh; it computes AUTS and the UE sends
        Authentication Failure (synch failure).
  \item The MME discards its remaining vectors for this subscriber and sends AIR with
        Re-Synchronization-Info = RAND $\Vert$ AUTS.
  \item The AuC verifies MAC-S, sets its SQN to the USIM's value and generates new vectors.
  \item The MME restarts authentication with a new vector; it now succeeds.
\end{steps}
\end{procedurebox}

\begin{acnfigure}
  \resizebox{\linewidth}{!}{\acnfig{lecture-11/sqn-resync}}
  \caption{Sequence-number re-synchronisation (\cref{proc:resync}).}
  \label{fig:resync}
\end{acnfigure}

% ---------------------------------------------------------------------
\section{Worked examples}

\begin{examplebox}{Size of an AIA}\label{ex:aia}
\textbf{Problem.} Compute the size of the Authentication-Info AVP carrying five E-UTRAN vectors (all AVPs
vendor-specific, optional Item-Number omitted).

\textbf{Solution.} Using the AVP length rule of Lecture~10 with a 12-byte header: RAND $12+16 = 28$, XRES (64
bits) $12+8 = 20$, AUTN $12+16 = 28$, $K_{\mathrm{ASME}}$ $12+32 = 44$. One E-UTRAN-Vector: $12 + 120 = 132$
bytes; five vectors in the grouped Authentication-Info: $12 + 5\times132 = 672$ bytes.
\end{examplebox}

\begin{examplebox}{Restore from backup}\label{ex:resync}
\textbf{Problem.} An HSS of $10^7$ subscribers is restored from a one-day-old backup, so its SQN counters
are behind the USIMs'. If all subscribers authenticate within the next hour, how much extra S6a load is
created?

\textbf{Solution.} Each subscriber's first authentication fails with a synchronisation failure and needs
one extra AIR: $10^7/3600 = 2778$ extra AIR per second for an hour -- on top of normal traffic and
comparable to it. Operators avoid this by replicating SQN updates continuously rather than relying on
backups, or by advancing restored counters by a safe margin.
\end{examplebox}

\begin{examplebox}{Choosing a redundancy scheme}\label{ex:geo}
\textbf{Problem.} Sites have $a = 0.999$. Compare one site, two sites each able to carry the full load,
and three sites of which any two can carry the load.

\textbf{Solution.} From \cref{eq:avail}: one site $u = 10^{-3}$ (8.76\,h/year); 1-of-2 $u^2 = 10^{-6}$
(31.5\,s/year); 2-of-3 $3u^2-2u^3 \approx 3.0\times10^{-6}$ (94.5\,s/year). The 2-of-3 design is three
times less available than 1-of-2 but needs only 50\,\% capacity per site instead of 100\,\%, i.e.\ 150\,\%
total instead of 200\,\%.
\end{examplebox}

% ---------------------------------------------------------------------
\section{Technology comparison}
\begin{table}[htbp]
\centering\small
\caption{Subscriber-data functions across generations.}\label{tab:cmp}
\rowcolors{2}{acnmist}{white}
\begin{tabularx}{\linewidth}{@{}l Y Y Y@{}}
\toprule
\rowcolor{white} & HLR/AuC (2G/3G) & HSS (4G, IMS) & UDM/AUSF/UDR (5G) \\
\midrule
Protocol         & MAP over SS7/SIGTRAN & Diameter (S6a, Cx, Sh) & HTTP/2 SBI (Nudm, Nausf, Nudr) \\
Authentication   & GSM triplets, UMTS quintets & EPS vectors with $K_{\mathrm{ASME}}$ & 5G-AKA, EAP-AKA$'$ \\
Identity privacy & TMSI only            & TMSI/GUTI only         & Concealed SUPI (SUCI) \\
Data model       & Per-node records     & Profile per access + IMS & Common UDR, NF front-ends \\
Discovery        & Global titles        & Realms, DRAs           & NRF \\
\bottomrule
\end{tabularx}
\end{table}

\paragraph{HSS versus visitor databases} The HSS is the master copy; the MME (like the VLR in 2G/3G)
caches what it needs for its subscribers. Caching keeps the HSS out of most procedures -- a TAU within the
same MME needs no S6a message -- at the price of keeping copies consistent with IDR/DSR/CLR.

% ---------------------------------------------------------------------
\section{Key takeaways}
\begin{takeawaybox}
\begin{itemize}
  \item HSS = HLR + AuC + IMS user data; the single subscriber authority, reached over S6a, Cx, Sh,
        SWx and MAP.
  \item The AuC keeps $K$ secret and issues vectors (RAND, XRES, AUTN, $K_{\mathrm{ASME}}$) built with
        $f_1$--$f_5$ (MILENAGE).
  \item SQN = SEQ $\Vert$ IND proves freshness; failures are repaired with AUTS re-synchronisation.
  \item The HSS records the serving MME, cancels old registrations and pushes profile changes.
  \item Storage is small ($\approx$100\,GB for $10^7$ users); availability and consistency are the hard
        parts, modelled by $k$-of-$N$ redundancy with independence as the key assumption.
  \item The front-end/repository split (UDC) leads directly to UDM/AUSF/UDR in 5G.
\end{itemize}
\end{takeawaybox}

% ---------------------------------------------------------------------
\section{Exercises}

\begin{exercisebox}{1}{Master and cache}\label{exr:split}
Which subscriber data are stored permanently in the HSS and which are cached in the MME? Which messages
keep the cache consistent?
\end{exercisebox}

\begin{exercisebox}{1}{Realm from IMSI}\label{exr:realm}
For IMSI 432101234567890 with a two-digit MNC, write the EPC home-network realm that a visited MME uses
as Destination-Realm.
\end{exercisebox}

\begin{exercisebox}{2}{Storage dimensioning}\label{exr:storage}
Compute the memory needed for $5\times10^7$ subscribers with a 6\,kB profile, three replicas and a fill
factor of 0.7.
\end{exercisebox}

\begin{exercisebox}{2}{How many vectors per request?}\label{exr:batch}
With $\lambda_{\mathrm{auth}} = 1157$\,s$^{-1}$, compute the AIR rate for $n = 1, 3, 5$. If a subscriber uses on
average two vectors per MME visit before moving to another MME, what fraction of generated vectors is
wasted for $n = 5$?
\end{exercisebox}

\begin{exercisebox}{3}{Redundancy design}\label{exr:kofn}
Derive the expressions for $A_{2\text{-of-}3}$ and $A_{1\text{-of-}2}$ from \cref{eq:avail}. Then suppose that,
besides independent failures, a common-mode failure (e.g.\ a faulty software release) disables all sites
with probability $p_c = 10^{-5}$. Recompute the 1-of-2 downtime and comment.
\end{exercisebox}

\begin{exercisebox}{3}{The IND array}\label{exr:ind}
MME~A holds an unused vector with IND $= 7$. The HSS keeps issuing vectors to MME~B, incrementing IND
cyclically. After how many newly issued vectors can the vector held by MME~A become unacceptable, and
under which condition does it actually fail?
\end{exercisebox}

\subsection*{Solutions}
\addcontentsline{toc}{subsection}{Solutions}

\begin{solutionbox}{\ref{exr:split}}
Permanent in the HSS: IMSI, MSISDN, $K$/OPc and SQN (AuC), the subscription profile and the identity of the
serving MME. Cached in the MME: the profile (from ULA), unused vectors and the security context. Consistency:
IDR/DSR (profile changes), CLR (withdrawal), PUR (MME deletes the subscriber), RSR (HSS restart).
\end{solutionbox}

\begin{solutionbox}{\ref{exr:realm}}
MCC $= 432$, MNC $= 10$ (padded to three digits as 010), so the realm is
\[\texttt{epc.mnc010.mcc432.3gppnetwork.org}\]
\cite{ts23003}. The MNC length is not encoded in the IMSI; the MME needs it from configuration.
\end{solutionbox}

\begin{solutionbox}{\ref{exr:storage}}
$M = 5\times10^7\times6\times10^3\times3/0.7 = 1.29\times10^{12}$\,bytes $\approx 1.3$\,TB -- distributed over
a cluster, still feasible in memory.
\end{solutionbox}

\begin{solutionbox}{\ref{exr:batch}}
$\lambda_{\mathrm{AIR}} = 1157,\ 386,\ 231$\,s$^{-1}$. With $n=5$ each visit triggers one AIR and uses two of five
vectors: three of five (60\,\%) are wasted, so the AuC computes $5/2 = 2.5$ times as many vectors as needed.
With $n = 3$, one AIR per visit and 33\,\% waste. The best $n$ matches the typical number of authentications
per visit.
\end{solutionbox}

\begin{solutionbox}{\ref{exr:kofn}}
$A_{1\text{-of-}2} = 2a(1-a) + a^2 = 1-(1-a)^2$. $A_{2\text{-of-}3} = 3a^2(1-a) + a^3$, so $1-A = 1 - 3a^2 + 2a^3
= 3u^2 - 2u^3$ with $u = 1-a$. With a common mode: $1-A \approx p_c + u^2 = 10^{-5} + 10^{-6} = 1.1\times10^{-5}$,
i.e.\ 347\,s per year instead of 31.5\,s. The common-mode term dominates by a factor of ten: beyond two sites,
staged software roll-outs and configuration discipline improve availability more than extra hardware.
\end{solutionbox}

\begin{solutionbox}{\ref{exr:ind}}
IND takes 32 values, so the next vector with IND $= 7$ (and a higher SEQ) is issued after 32 new vectors.
The vector held by MME~A fails only if that newer vector is \emph{used} by the USIM before MME~A uses its own:
the USIM then stores the higher SEQ in slot 7 and rejects the older one. Up to 31 newer vectors can therefore
be used in any order without affecting it.
\end{solutionbox}

% ---------------------------------------------------------------------
\section{Further reading}
The HSS is defined in TS~23.002 \cite{ts23002}; identities and realms in TS~23.003 \cite{ts23003}. The
authentication framework and the SQN management annex are in TS~33.102 \cite{ts33102}, the EPS-specific parts in
TS~33.401 \cite{ts33401}, and MILENAGE in TS~35.206 \cite{ts35206}. S6a is specified in TS~29.272 \cite{ts29272},
MAP in TS~29.002 \cite{ts29002}. Olsson et al.\ \cite{olsson2013epc} describe the HSS from a deployment
perspective.

\printbibliography[heading=subbibliography,title={\ACNstrRefs}]
\end{document}
